\begin{proof}
\begin{enumerate}
\item \textbf{Degenerate cases.}  
If \(x=0\) then \(\|x+y\|=\|y\|=\|0\|+\|y\|\); similarly if \(y=0\). Hence the inequality holds with equality.  
Henceforth assume \(x\neq0\) and \(y\neq0\).

\item \textbf{Expand \(\|x+y\|^{2}\).}  
By definition of the norm and bilinearity of the inner product,
\begin{align*}
\|x+y\|^{2}
   &=\langle x+y,x+y\rangle                                      \\
   &=\langle x,x\rangle+\langle x,y\rangle
     +\langle y,x\rangle+\langle y,y\rangle .
\end{align*}
Since \(\langle y,x\rangle=\overline{\langle x,y\rangle}\),
\begin{align}
\|x+y\|^{2}
   &=\|x\|^{2}+ \|y\|^{2}
     +\langle x,y\rangle+\overline{\langle x,y\rangle} \nonumber\\
   &=\|x\|^{2}+ \|y\|^{2}+2\Re\langle x,y\rangle .
\tag{1}
\end{align}

\item \textbf{Estimate the middle term.}  
For any complex number \(z\),
\[
|\Re z|\le |z| .
\tag{2}
\]
Indeed, \(|z|^{2}=(\Re z)^{2}+(\Im z)^{2}\ge (\Re z)^{2}\), whence \(|\Re z|\le|z|\).  
Applying (2) with \(z=\langle x,y\rangle\) gives
\[
\Re\langle x,y\rangle\le |\langle x,y\rangle| .
\tag{3}
\]

Now invoke the Cauchy–Schwarz inequality:
\[
|\langle x,y\rangle|\le \|x\|\,\|y\| .
\tag{4}
\]

Combining (3) and (4) we obtain
\[
\Re\langle x,y\rangle\le \|x\|\,\|y\| .
\tag{5}
\]

\item \textbf{Upper bound for \(\|x+y\|^{2}\).}  
Insert (5) into (1):
\begin{align}
\|x+y\|^{2}
   &\le \|x\|^{2}+2\|x\|\,\|y\|+\|y\|^{2} \nonumber\\
   &=\bigl(\|x\|+\|y\|\bigr)^{2}.
\tag{6}
\end{align}
All quantities involved are non‑negative.

\item \textbf{Take square roots.}  
Since the square‑root function is monotone on \([0,\infty)\),
\[
\|x+y\|\le \|x\|+\|y\|.
\]

\item \textbf{Equality condition.}  
Equality in the chain (6) occurs precisely when each inequality used becomes an equality.

\begin{itemize}
\item From (5): equality in \(\Re\langle x,y\rangle\le |\langle x,y\rangle|\) forces \(\langle x,y\rangle\) to be a non‑negative real number (i.e. \(\Im\langle x,y\rangle=0\) and \(\Re\langle x,y\rangle\ge0\)).
\item Equality in Cauchy–Schwarz (4) holds iff \(x\) and \(y\) are linearly dependent, i.e. \(y=\lambda x\) for some \(\lambda\in\mathbb{C}\).
\item Substituting \(y=\lambda x\) gives \(\langle x,y\rangle=\lambda\langle x,x\rangle\); this is real and non‑negative exactly when \(\lambda\) is a non‑negative real scalar.
\end{itemize}
Thus equality \(\|x+y\|=\|x\|+\|y\|\) holds iff one of the following is true:

\begin{enumerate}
\item \(x=0\) or \(y=0\) (the degenerate case);
\item There exists \(\lambda\in\mathbb{R}\) with \(\lambda\ge0\) such that \(y=\lambda x\) (the vectors are positively colinear).
\end{enumerate}
These are precisely the cases announced in the statement of the triangle inequality.
\end{enumerate}
\end{proof}